\section{Additional simulation results}\label{app:sim}
% =============================================================================================
% Owner: simulation writer. Tables behind Section 5: tab:e2-size, tab:e2dep-size,
% tab:e2-power-ci, tab:e2dep-power, tab:e2dep-check (manual, paper/scripts/e2dep_check.py),
% fig:e2-llr-scatter, tab:e1-size-sim-full (generated e1_size_sim.tex), tab:e1-sim-dependence and
% tab:e1-sim-z (manual, paper/scripts/e1_sim_diag.py), tab:e1-fixedb-check (manual,
% paper/scripts/e1_fixedb_check.py). New labels: app:sim-size, app:sim-power,
% app:sim-dep, app:sim-llr, app:sim-e1, tab:e2dep-check, tab:e1-size-sim-full,
% tab:e1-sim-dependence, tab:e1-sim-z, tab:e1-fixedb-check. Readable versions of e2_power_ci / e2dep_power:
% paper/scripts/e2_power_ci_app.py (cell-by-cell identical to the generated tables).
% =============================================================================================

This appendix collects the tables behind \cref{sec:sim}: the size of every test in E2 and E2-dep,
the power with Wilson 95\% intervals, the marginal check of the dependence-only families, the
learned against the exact log-likelihood ratios, and the full results and diagnostics of the
simulated panel placebo E1(a). Designs and replication counts are those of \cref{sec:sim-design}.
The generated tables and figures write $\delta$ for the deviation parameter $\vartheta$,
$\alpha+\beta$ for the GARCH persistence $a+b$, and $b$ for the E2-dep block length $L_b$.

% ---------------------------------------------------------------------------------------------
\subsection{Size with independent windows}\label{app:sim-size}

\Cref{tab:e2-size,tab:e2dep-size} give the rejection rates under $H_0$ at the three levels. Every
window test uses a permutation p-value computed on independent windows, which is exact
(\cref{prop:exact-level}), so these tables check the implementation: at 5\% every interval contains
the nominal level, and over the three levels 3 of the 150 intervals of the window tests miss it,
all at 1\% or 10\% (\cref{sec:sim-size}). KS is marked with a dagger where returns within a window
are dependent; its p-value is then invalid. At 5\% its interval excludes 0.05 in family (c) with
$n=500$ and in the E2-dep mixture family; in family (c) with $n=2000$ the lower bound of its interval
is 0.050.

\begin{table}[p]
  \centering
  \caption{Size of the E2 tests: rejection rates under $H_0$ ($Q=P$) with Wilson 95\% intervals,
    independent windows of $L=20$ days, $n$ windows per sample, $N=200$ replications per cell, at
    $\alpha\in\{0.01,0.05,0.10\}$.}
  \label{tab:e2-size}
  \scriptsize
  \fittable{e2_size}
  \tablenote{\inputnote{e2_size_note}}
\end{table}

\begin{table}[tbp]
  \centering
  \caption{Size of the E2-dep tests: rejection rates under $H_0$ ($\vartheta=0$ in the mixture
    family, $L_b=20$ in the block family, independent seeds) with Wilson 95\% intervals,
    independent windows of $L=20$ days, $n=2000$ windows per sample, $N=100$ replications per
    cell, at $\alpha\in\{0.01,0.05,0.10\}$.}
  \label{tab:e2dep-size}
  \small
  \fittable{e2dep_size}
  \tablenote{\inputnote{e2dep_size_note}}
\end{table}

% ---------------------------------------------------------------------------------------------
\subsection{Power with confidence intervals}\label{app:sim-power}

\Cref{tab:e2-power-ci,tab:e2dep-power} repeat the power of \cref{fig:e2-power,fig:e2dep-power}
with Wilson 95\% intervals. With 100 replications per alternative in E2 and 50 in E2-dep, the
intervals are wide (for instance 0.46 [0.33, 0.60] for the \WC{} test at $\vartheta=0.5$ in the
mixture family), so only large differences between window tests are resolved; the rankings stated
in \cref{sec:sim-power} compare point estimates. The differences between KS and the window tests
in families (a), (b) and (d), and between the window tests and KS in E2-dep, are far larger than
the intervals.

\begin{table}[p]
  \centering
  \caption{Power against the four E2 families with Wilson 95\% intervals: rejection rates at
    $\alpha=0.05$, independent windows of $L=20$ days, $n$ windows per sample, 200 replications at
    $k=0$ ($H_0$) and 100 at each alternative $k=1,\dots,4$; the levels of each family are listed
    in its heading.}
  \label{tab:e2-power-ci}
  \scriptsize
  \fittable{manual_e2_power_ci_app}
  \tablenote{p-values as in \cref{tab:e2-size} (LR oracle: Monte Carlo, 2000 null samples).
    $^{\dagger}$KS is not valid in family (c); its power there is not size-adjusted.
    $^{\ddagger}$Approximate oracle ($M=512$ Monte Carlo draws of the start variance).}
\end{table}

\begin{table}[tbp]
  \centering
  \caption{Power against the dependence-only families of E2-dep with Wilson 95\% intervals:
    rejection rates at $\alpha=0.05$, independent windows of $L=20$ days, $n=2000$ windows per
    sample, 100 replications under $H_0$ ($k=0$) and 50 at each alternative.}
  \label{tab:e2dep-power}
  \small
  \fittable{manual_e2dep_power_app}
  \tablenote{Family (a): a share $\vartheta$ of the $Q$ windows consists of $L$ independent
    returns; family (b): $Q$ windows are concatenations of independent $P$ blocks of length
    $L_b$. Every return of $Q$ has exactly the marginal law of $P$ (check in
    \cref{tab:e2dep-check}). p-values as in \cref{tab:e2-power-ci}; the oracle uses 1000 Monte
    Carlo null values. $^{\dagger}$KS sees only the marginal and has no power by construction; the
    returns of a GARCH window share their conditional variance, so its p-value is not valid either
    (rates under $H_0$ in \cref{tab:e2dep-size}); its rejection rates are size distortions, not
    power. The oracle is approximate: the window and block densities of $P$ integrate the
    conditional variance at the window or block start over 512 draws, the position marginals over
    2048 draws; statistic and null use the same draws, so the oracle test has exact level but is
    only approximately most powerful.}
\end{table}

% ---------------------------------------------------------------------------------------------
\subsection{The dependence-only families: marginal check and diagnostics}\label{app:sim-dep}

\Cref{tab:e2dep-check} shows what changes along the two E2-dep families and what does not. Within a
window, the lag-1 autocorrelation of $\abs{r-\mu}$ falls from 0.231 under the null to 0.002 at
$\vartheta=1$ in the mixture family, and from 0.223 to 0.000 between $L_b=20$ and $L_b=1$ in the
block family. The standard deviation of the log realised variance across windows falls from 0.846
to 0.585 and from 0.840 to 0.591: windows with independent returns have less dispersed realised
volatility. The variance of the returns stays between $3.91\times10^{-4}$ and $4.19\times10^{-4}$
at every level; this \GARCHt{} law has no finite fourth moment, so sample variances over 20,000
windows vary by a few per cent. Kurtosis is not reported for the same reason.

The marginal check compares, at the \iid{}-marginal level of each family, $10^6$ returns of $P$
(one per independent window, so that they are independent) with $10^6$ returns of $Q$. For the
block family the two-sample KS statistic is $D=0.0011$ with $p=0.616$. For the mixture family it is
$D=0.0027$ with $p=0.001$. The two extreme levels are the same law by construction (the mixture
generator at $\vartheta=1$ calls the generator of $L_b=1$), and the two checks use the same $P$
returns and independent $Q$ seeds. We therefore read the mixture result as a chance rejection of
the check, not as a difference between marginals, but we have not investigated it further. A
discrepancy of this size is in any case far below what the experiment can detect: in the E2-dep
runs, KS rejects at $\vartheta=1$ at rate 0.10, the same as under $H_0$ (\cref{tab:e2dep-power}).

\begin{table}[tbp]
  \centering
  \caption{E2-dep dependence diagnostics and marginal check. Columns 3--5: within-window lag-1
    autocorrelation of $\abs{r-\mu}$, cross-window standard deviation of the log realised
    variance, and return variance, each on 20,000 simulated $Q$ windows of $L=20$ days per level.
    Columns 6--7: two-sample KS test of the marginal at the \iid{}-marginal level, $10^6$ returns
    of $P$ (one per independent window) against $10^6$ returns of $Q$.}
  \label{tab:e2dep-check}
  \small
  \fittable{manual_e2dep_check}
\end{table}

% ---------------------------------------------------------------------------------------------
\subsection{Learned versus exact log-likelihood ratios}\label{app:sim-llr}

\Cref{fig:e2-llr-scatter} plots the learned logit against the exact $\log(p/q)$ for one replication
of E2 part f. With early stopping (top row) the logit is nearly constant: the early-stopped
classifier separates the two laws barely better than chance. Without early stopping (bottom row)
it is overconfident and scattered, with little relation to the exact log ratio (root mean squared
errors 4.917 and 4.682 in \cref{tab:e2-llr}). For GBM against Merton the exact log ratio is close to zero for most
windows and large only for the few Merton windows with a big jump; no learned score picks these
windows out, which is why every window test has near-zero power in family (d)
(\cref{sec:sim-llr}).

\begin{figure}[tbp]
  \centering
  \inputfigure{e2_llr_scatter}
  \caption{Learned logit $\lhat(w)$ against the exact log-likelihood ratio $\log p(w)/q(w)$ on 2000
    fresh windows of $L=20$ days (1000 per law), replication 0 of E2 part f; logits fitted on 2000
    windows per law. Left: GBM against Merton; right: GARCH persistence 0.98 against 0.90 (window
    density integrated over the start variance by Monte Carlo). Top: test default (early stopping,
    at most 1000 steps, learning rate $10^{-3}$); bottom: 1000 steps without early stopping.
    Dashed: identity. The correlations and slopes printed in the panels are those of this
    replication; medians over 20 replications are in \cref{tab:e2-llr}.}
  \label{fig:e2-llr-scatter}
\end{figure}

% ---------------------------------------------------------------------------------------------
\subsection{The simulated panel placebo}\label{app:sim-e1}

\paragraph{Calibration of the placebo panel.}
\Cref{tab:e1-sim-dependence} compares the cross-sectional dependence of the simulated panels with
that of the real panel. Against real \TEST{}, whose dates the absolute tests use, $\Rnew^*$ has a
somewhat higher mean pairwise return correlation (median over replications 0.319, interquartile
range 0.267 to 0.407, against 0.254) and a similar mean pairwise correlation of the rolling 20-day
log-volatility (median 0.270 against 0.309). Against real \TRAIN{}, which enters the relative
tests, $\Rpast^*$ has much less volatility co-movement (median 0.327 against 0.611). The factor
GARCH-$t$ model has a single common volatility process and residual volatilities that are
independent across tickers, so it understates the common volatility regimes of 2018--2022. In the
cell indep both correlations of $G^*$ are zero in population, by construction.

\begin{table}[tbp]
  \centering
  \caption{Cross-sectional dependence of the real panel and of the E1(a) placebo panels (100
    tickers): mean pairwise correlation of daily returns and of the rolling 20-day log-volatility.
    Placebo rows: median over 200 replications of these means, interquartile range in brackets.}
  \label{tab:e1-sim-dependence}
  \small
  \fittable{manual_e1_sim_dependence}
\end{table}

\paragraph{Fixed-$b$ references.}
The fixed-$b$ p-values of \cref{tab:e1-size-sim} were computed after the run from the stored
studentised statistics, the numbers of evaluation dates $U_X,U_Y$ and the Bartlett lags $H_X,H_Y$
of the two samples \citep{kiefer_2005_new}. The conservative reference is that of a single series
with $\min(U_X,U_Y)$ evaluation dates and lag $\max(H_X,H_Y)$. We expect it to be the
heaviest-tailed case of the two-sample statistic whatever the share $\lambda\in[0,1]$ of the first
sample in the variance of the mean difference, but we have not proved this, and with unequal lags
it is not obvious. \Cref{tab:e1-fixedb-check} checks it by simulation for the designs of E1. For
each of the 154 combinations $(U_X,H_X,U_Y,H_Y)$ that occur in the fixed-$b$ rows (65 for the
absolute tests, 89 for the anchored test), 200,000 pairs of contiguous \iid{} Gaussian series give
the two-sample statistic at 21 values of $\lambda$ on $[0,1]$, and each value is referred to the
conservative reference with the reference draws of the run. For $0<\lambda<1$ the largest
rejection rate over all combinations is 0.0094, 0.0501 and 0.1006 at the 1\%, 5\% and 10\% levels,
below the nominal level or within Monte Carlo error of it. The largest rates over the whole grid,
0.0110, 0.0514 and 0.1018, are attained at $\lambda\in\{0,1\}$ in combinations where the statistic
is then a single series with the same number of dates and the same lag as the reference, so they
measure the Monte Carlo error of the 40,000 reference draws. The check supports the conservative
reference for these designs; it is not a proof for other numbers of dates or lags. The
equal-variance reference is the two-sample limit with each sample's own number of dates and lag
and equal per-date variances. Both references, and the check, treat the evaluation dates as
contiguous, whereas the Driscoll--Kraay estimator \eqref{eq:dk-var} runs on the date grid with the
training blocks as gaps, so they are approximations.

\begin{table}[tbp]
  \centering
  \caption{Monte Carlo check of the conservative fixed-$b$ reference on the E1(a) designs. For
    every combination of numbers of evaluation dates $U_X,U_Y$ and Bartlett lags $H_X,H_Y$ that
    occurs in the fixed-$b$ rows of \cref{tab:e1-size-sim-full} (Calls: test calls with these
    combinations, over tests, cells and 200 replications), 200,000 pairs of contiguous \iid{}
    $N(0,1)$ series of lengths $U_X$ and $U_Y$ give the two-sample Bartlett HAC statistic with
    share $\lambda$ of the first sample in the variance of the mean difference, for 21 values of
    $\lambda$ on $[0,1]$. Each value is referred to the conservative reference (one series with
    $\min(U_X,U_Y)$ dates and lag $\max(H_X,H_Y)$, same reference draws as in the run). Entries:
    largest rejection rate over combinations and $0<\lambda<1$; in parentheses, over
    $0\le\lambda\le1$ ($\lambda\in\{0,1\}$: a single series).}
  \label{tab:e1-fixedb-check}
  \small
  \fittable{manual_e1_fixedb_check}
  \tablenote{Monte Carlo standard error of one rate, $\sqrt{\alpha(1-\alpha)/200{,}000}$: 0.0002,
    0.0005 and 0.0007 at the three levels, to which the error of the 40,000 reference draws adds.
    Script: \texttt{paper/scripts/e1\_fixedb\_check.py}.}
\end{table}

\paragraph{The studentised statistic under the null.}
\Cref{tab:e1-sim-z} summarises $\zhac$ over the placebo replications. If the Driscoll--Kraay
p-value were calibrated, $\zhac$ would be approximately standard normal. In the joint cell its mean
is close to zero (between $-0.04$ and 0.12) and its standard deviation between 1.12 and 1.19, the
mild over-dispersion that 15 evaluation dates and a median lag of one window produce
(\cref{rem:small-T}). In the indep cell the \WC{} statistic has mean 0.74 and standard deviation
1.76, and the \LRC{}, anchored and two-critic statistics have means between 0.33 and 0.39: the
distribution is shifted, not only spread. The C2ST statistic is over-dispersed (standard deviation
1.42) but not shifted (mean 0.03). The paired comparison of the two cells, which share their
real panels and test seeds, gives exact McNemar p-values \citep{mcnemar_1947_note} of
$2.0\times10^{-6}$ for the \WC{} test, 0.013 for the \LRC{} test, 0.017 for the anchored test and
$4.2\times10^{-5}$ for the two-critic test (Driscoll--Kraay p-values at 5\%), but 0.17 for C2ST and
0.26 and 0.29 for the wild-bootstrap MMD and \SigMMD{} tests. In the joint cell all protocol
p-values reject between 0.045 and 0.090 at 5\%; in the indep cell the wild-bootstrap kernel tests
are the closest to nominal (0.075 for MMD and 0.100 for \SigMMD{}, against 0.115 to 0.270 for the
learned and relative tests; \cref{tab:e1-size-sim-full}). \Cref{tab:e1-sim-se} in
\cref{sec:sim-dependent} and \cref{tab:manual-e1-sim-stat} in \cref{sec:theory-dependence}
separate the statistic from its standard error: the shift of $\zhac$ in the indep cell comes with
a mean of $T$ close to that of the joint cell and a negative correlation between $T$ and its
standard error.

\begin{table}[tbp]
  \centering
  \caption{The Driscoll--Kraay statistic $\zhac$ under the E1(a) placebo null (time-block split),
    200 replications per cell: mean and standard deviation over replications, share above the
    normal 5\% critical value (the size in \cref{tab:e1-size-sim}), and the medians of the number
    of evaluation dates and of the Bartlett lag $H$ of the first sample ($\Rnew^*$ for the
    absolute tests, $\Rpast^*$ for the anchored test; not recorded for the two-critic test).}
  \label{tab:e1-sim-z}
  \small
  \fittable{manual_e1_sim_z}
\end{table}

\begin{table}[p]
  \centering
  \caption{Size of the real-data pipeline in the simulated factor panel (E1(a)), all p-value
    types: rejection rates under $H_0$ with Wilson 95\% intervals at
    $\alpha\in\{0.01,0.05,0.10\}$, 200 replications per cell, 100 tickers, $L=20$, 3000 real and
    3000 synthetic \TEST{} windows, 6200 past windows. $G^*$ joint: independent draw of the whole
    factor panel; indep.: tickers drawn independently with the same marginal laws.}
  \label{tab:e1-size-sim-full}
  \scriptsize
  \fittable{e1_size_sim}
  \tablenote{Learned and relative tests use the time-block split unless ``ticker split'' is
    stated. $^{\star}$Protocol p-value (Driscoll--Kraay ``HAC'' for learned and relative tests,
    wild bootstrap for kernel tests). Fixed-$b$ rows: the same statistic against the conservative
    and the equal-variance fixed-$b$ references (\cref{app:sim-e1}). ``permutation'': for learned
    and relative tests, permutation of the evaluation labels of the critic trained on time blocks;
    for MMD and \SigMMD{}, \iid{} permutation of all windows. ``stud.\ permutation'': studentised
    permutation of the anchored statistic. KS: asymptotic p-value, which assumes \iid{} returns and
    is invalid here.}
\end{table}

\clearpage
